% !TeX root = ../papers/nc-p3.tex
\section{十五次元閉包の有限四辺での全正則性判定}
\label{sec:main-regular-specialization}
\begin{theorem}\label{thm:main-regular-specialization}
十五次元成分の閉包で有限四辺の標準有限Frobenius候補を取る。
その左右AS正則かつKoszulであるための必要十分条件は、
経路長4の四中心元で命題\ref{prop:central-arrow-module-cutoff}の検査が通ることである。
経路次元を14次、中心剰余を13次、左右標準複体を16次まで調べればよい。
この閉包の有限四辺の正則軌跡は開であり、その全点で合成整域性・中心有限性・頂点指数4が成立する。
各点の添字保持同型は偏極中心と零番隅のKummer印で決まる。
\end{theorem}
\begin{proof}
正則点を取り、定理\ref{thm:finite-short-trace-pole}の離散付値環 $O$ 上の特殊化を使う。
$t$ を一様化元、$\mathcal B$ を経路格子、$C$ を一般の行列traceが生成する $O$ 代数、
$T=C\mathcal B$ と置く。$T$ は有限型可換環 $C$ の有限加群で左右Noetherである。
一般・特殊とも正則なので各次数の経路格子は自由であり、$\mathcal B$ は $T$ の部分環である。
下添字 $s$ を $t$ による剰余とする。ここでは $B_s\to T_s$ の単射性を必要としない。

$C_s$ の原点以外の閉指標で $T_s$ の繊維のどの単純も頂点単純ではない。
実際、一般絶対単純の特殊化の半単純型にはこの繊維の全単純型が現れる
（特殊化traceとCayley--Hamilton恒等式、定理\ref{thm:pi-vertex-extension-criterion}の証明）。
全頂点の次元は4である。全因子が頂点単純なら全正次数traceが零になり原点となる。
頂点と非頂点が混在すると仮定する。定理\ref{thm:four-matrix-balanced-simples}により
非頂点は均衡し、非薄なものは重複込みで高々一つで、その次元は $2\delta$ または $3\delta$ である。
薄いものと頂点の第一Extは両向きとも零、この一つと全頂点の第一Extも少なくとも一向きは全て零である。

この片向きの消滅は、関係の係数も変わる相対族で一般絶対単純への平滑化を妨げる。
特殊格子表現を半単純化してもその点は絶対単純軌跡の閉包にある。実際格子表現へ特殊化してから
特殊繊維の頂点基底変更の一径数群で組成列の随伴次数表現へ移せる。
この閉包の中で半単純点に至る曲線を取り直し、完備離散付値環を使う。
二群に分けた矢行列の指定側のブロックを $t^m$ まで消したとする。
関係のそのブロックの $t^m$ 係数は特殊対角ブロックに関する一次コサイクルである。
指定側を全て零にすれば残る行列は三角で、その側は変形された関係のどの積でも恒等的に零だからである。
消滅方向の第一Extにより頂点ごとの基底変更で次の係数も消える。
完備極限で一般繊維に非自明な不変部分を得て矛盾する。

従って $T_s/T_s B_{s,+}$ は有限次元である。この商は $C_s$ 上有限で、非原点に台があれば
繊維商の単純商は正次数経路を消し、上で除いた頂点型となる。
ここで正次数イデアルが作る部分加群は $T_s=C_s B_s$ より両側イデアルである。
次数帰納により $T_s$ は $B_s$ 上右有限で、生成元には $C_s$ の有限個の斉次元を選べる。
実際各次数の元は $C_s$ の元と $B_s$ の経路の積の和であり、
正次数経路を除いた剰余では中心係数と頂点冪等元だけで張られる。

その中心係数を $c_1,\ldots,c_r\in C$ へ斉次に持ち上げると
\[
                      T=\sum_{j=1}^r c_j\mathcal B .
\]
各次数で右辺を割った余核は有限 $O$ 加群で、特殊化して零だからNakayamaで零である。
全次数で同じ有限個の $c_j$ を使っている。$c_j$ は中心にあるから左加群でも同じ等式になる。
有限中心化拡大のNoether降下定理を $\mathcal B\subset T$ に使えば
$\mathcal B$ は左右Noether、従って $B_s=\mathcal B/t\mathcal B$ も左右Noetherである。
ここで降下を使う包含は付値環上にあり、剰余写像の核については何も仮定していない。
定理\ref{thm:noetherian-pi-window-cutoff}から合成整域、
定理\ref{thm:finite-domain-specialization-parameters}から指数4と中心のHilbert関数の保存を得る。
定理\ref{thm:main-domain-quartic-parameters}が有限検査の両含意と同型判定を与える。

最後に被約な成分の中で各正則点の近傍を考える。必要な有限次数の経路空間はそこで
一般点と同次元なので、小行列を選べばベクトル束になる。標準有限Frobeniusデータも同じ階数開条件で取る。
四次中心核も一般点と同じ五次元で、交換子写像の一定階数開集合上の部分束である。
ある四元の検査が通ればその局所切断を持ち上げられる。中心剰余の次数別次元の所定上界は開条件であり、
自由な経路項を持つ左右の各有限標準複体の指定位置での完全性も開条件である。
中心剰余の次元に等号を使う場合も、この近傍では一般点の下界を使える。
実際特殊点で正規化を与える四元は、その固定した有限次数の剰余上界が一般点でも保たれるので
一般点でも有限剰余を与える。一般点の正則Cohen--Macaulay窓では四元が正則列となる。
これで剰余の所定次元が一般点で成立する。有限四辺の条件も射影的な辺の繊維次元の上半連続性による開条件である。
\end{proof}
